\documentclass[11pt]{article}
\usepackage[margin=0.85in]{geometry}
\usepackage{lmodern}
\usepackage{microtype}
\usepackage{enumitem}
\usepackage{listings}
\usepackage{needspace}
\usepackage[T1]{fontenc}

\lstset{
  basicstyle=\ttfamily\small,
  columns=fullflexible,
  keepspaces=true,
  showstringspaces=false,
  breaklines=true,
  frame=single,
  xleftmargin=0.4em,
  xrightmargin=0.4em,
  aboveskip=0.6em,
  belowskip=0.6em
}

\pagestyle{empty}
\setlength{\parskip}{0.4em}
\setlength{\parindent}{0pt}
\setlist[enumerate,1]{leftmargin=1.4em, itemsep=0.9em, topsep=0.4em}
\setlist[enumerate,2]{leftmargin=1.5em, itemsep=0.45em, topsep=0.3em, label=\arabic*.}
\setlist[itemize]{leftmargin=1.3em, itemsep=0.2em, topsep=0.25em}

\begin{document}

\begin{center}
{\LARGE\bfseries Midterm 2 --- Sample 1}\\[0.35em]
DATA 412/612
\end{center}

\vspace{0.3em}

Practice exam, questions only. Put your answers in an R Markdown document,
and make sure code is visible when you knit
(\texttt{knitr::opts\_chunk\$set(echo = TRUE)} in the first chunk).
Use the tidyverse rather than base R where you can.
You may use your notes, the R help files, and the internet.
Do not discuss the exam with anyone else.

\begin{enumerate}
\item \textbf{NBA} (8 pts).
The data frames below are a small slice of 2019 NBA statistics.

\texttt{teams}: \texttt{Team} (name), \texttt{W} (wins), \texttt{L} (losses),
\texttt{PSG} (points scored per game), \texttt{PAG} (points allowed per game).

\texttt{players}: \texttt{Name}, \texttt{teamID}, \texttt{Age},
\texttt{G} (games played), \texttt{GS} (games started),
\texttt{MP} (minutes per game), \texttt{FG} (field goals per game),
\texttt{FGA} (field-goal attempts per game).

\texttt{games}: \texttt{date\_time} (date and tip-off, stored as text),
\texttt{location}, \texttt{bulls} (Bulls points), \texttt{cavs} (Cavaliers points).

\begin{lstlisting}
library(tidyverse)

teams <- tribble(
  ~Team, ~W, ~L, ~PSG, ~PAG,
  "Cavaliers", 17, 53, 104, 113,
  "Bulls",     19, 52, 105, 113
)

players <- tribble(
  ~Name, ~teamID, ~Age, ~G, ~GS, ~MP, ~FG, ~FGA,
  "Zach LaVine",      "Bulls",     23, 62, 61, 34.5, 8.5, 18.1,
  "Lauri Markkanen",  "Bulls",     21, 48, 47, 32.7, 6.6, 15.2,
  "Otto Porter",      "Bulls",     25, 15, 15, 32.8, 6.5, 13.4,
  "Kevin Love",       "Cavaliers", 30, 16, 16, 27.4, 5.0, 13.0,
  "Jordan Clarkson",  "Cavaliers", 26, 69,  0, 27.1, 6.5, 14.7,
  "Collin Sexton",    "Cavaliers", 20, 70, 60, 31.0, 6.0, 14.3
)

games <- tribble(
  ~date_time, ~location, ~bulls, ~cavs,
  "Nov 10, 2018 8:00pm", "Bulls",      99,  98,
  "Dec 23, 2018 6:00pm", "Cavaliers", 112,  92,
  "Jan 21, 2019 1:00pm", "Cavaliers", 104,  88,
  "Jan 27, 2019 3:30pm", "Bulls",     101, 104
)
\end{lstlisting}

\begin{enumerate}
\item (4 pts) Use \texttt{parse\_datetime()} from \texttt{readr} to parse
\texttt{date\_time} in \texttt{games}, and assign the result back to
\texttt{games}.

\item (4 pts) Add the columns from \texttt{teams} onto \texttt{players}.
\end{enumerate}

\item \textbf{Strings} (10 pts).
Use only functions from the tidyverse. Your code should work for strings
of this shape in general, not only for the rows below.

\Needspace{14\baselineskip}
\begin{lstlisting}
dfdat <- tribble(
  ~frac, ~words,
  "4/5",  "Not",
  "7/8",  "all",
  "2/3",  "those",
  "4/7",  "who",
  "2/11", "wander",
  "8/13", "are",
  "4/6",  "lost."
)
\end{lstlisting}

\begin{enumerate}
\item (4 pts) Convert \texttt{frac} from a character vector of fractions
to a numeric vector of decimal values. For example, \texttt{"4/5"} should
become \texttt{0.8}. No \texttt{stringr} functions and no regular expressions.

\item (4 pts) Use a single regular expression to find every word in
\texttt{words} that begins with \texttt{wh} or \texttt{wa}, and replace
that prefix with \texttt{te}. So \texttt{who} becomes \texttt{teo}, and
\texttt{wander} becomes \texttt{tender}.

\item (2 pts) \textbf{Challenge.} Pig Latin moves the first consonant
cluster (consonants up to the first vowel) to the end of the word, then
appends \texttt{ay}. Examples: \texttt{hello} becomes \texttt{ellohay},
\texttt{study} becomes \texttt{udystay}, and \texttt{aloe} becomes
\texttt{aloeay}. Build the Pig Latin version of each string in
\texttt{words}.

For a little extra credit: keep an initial capital as a capital on the
new first letter (\texttt{John} becomes \texttt{Ohnjay}, not
\texttt{ohnJay}), and keep a trailing period in place
(\texttt{goodbye.} becomes \texttt{oodbyegay.}).
\end{enumerate}

\item \textbf{Tidying simulation output} (22 pts).
The file \texttt{mse\_mat.csv} is on the course page under Data.
It is a small extract of simulation output. Each row is one iteration.
The first six columns are the conditions:
\texttt{current\_seed}, \texttt{nullpi}, \texttt{Nsamp},
\texttt{ncontrols}, \texttt{cor}, and \texttt{poisthin}.

\texttt{cor} is a string. It holds either one number in parentheses,
such as \texttt{(0.5)}, or two numbers in parentheses separated by a
comma and a space, such as \texttt{(0.25, 0.25)}.

Every other column is a mean squared error for one combination of a
mean method and a variance method. The nine mean methods are
\texttt{ols}, \texttt{ruv2}, \texttt{ruv3}, \texttt{ruv4}, \texttt{cate},
\texttt{ruvbn}, \texttt{ruvbnl}, \texttt{ruvbnn}, and \texttt{ruvb}.
The variance methods are \texttt{o}, \texttt{m}, \texttt{c}, \texttt{l},
\texttt{lm}, and \texttt{lc}, plus a few longer labels that appear on
\texttt{ruv3} and \texttt{cate}. You do not need to know what the methods do.
Not every mean method has a variance method. The two pieces of the name
are separated by an underscore. For example, \texttt{ols\_m} is the
\texttt{ols} mean method with the \texttt{m} variance method, and
\texttt{ruvb\_} is the \texttt{ruvb} mean method with no variance method.

For each part, print the first six rows with \texttt{head()}.
The only non-tidyverse functions you may use in this question are
\texttt{head()} and \texttt{c()}.

\begin{enumerate}
\item (2 pts) Use \texttt{readr} to load the file, and view the first six rows.

\item (4 pts) Reshape the data so that, besides the six condition columns,
one column holds the mean/variance method name and another column holds
the mse for that method on that iteration.

\item (4 pts) From that method name, make two columns: the mean method and
the variance method. For \texttt{ruvbn}, \texttt{ruvbnl}, \texttt{ruvbnn},
and \texttt{ruvb}, the variance method should be the empty string
\texttt{""}.

\item (6 pts) Clean \texttt{cor}:
\begin{itemize}
  \item Remove every parenthesis.
  \item Replace the first number with that same number followed by an underscore.
  \item Remove every comma and the space after it.
\end{itemize}
For example, \texttt{(0.25, 0.25)} becomes \texttt{0.25\_0.25}, and
\texttt{(0.25)} becomes \texttt{0.25\_}.

\item (4 pts) Using the cleaned \texttt{cor}, make two numeric columns:
the first correlation and the second correlation. If \texttt{cor} had
only one value, the second correlation is \texttt{NA}.

\item (2 pts) Make an appropriate plot of how \texttt{Nsamp} and the mean
method relate to mean squared error.
\end{enumerate}
\end{enumerate}

\end{document}
